\section{Does failure survive additional alternatives?}\label{sec:expansion}
Let $X$ be an old set of complete plans, $Y$ a set of additions, and
$Z=X\cup Y$. Use the same supply function and compact projected-image
assumptions as before, with attained hull-dual prices where needed. Write
$D_A,\CH_A,V_A$ for the values associated with a set $A$. Suppose
$\Delta_X>0$. The following distinction is essential: a retained old plan
cannot be rescued by more deviations, but a newly admitted plan can become
a supported optimum.

\begin{proposition}[Expansion and an old dual certificate]\label{prop:expansion}
For every old $x$ and unchanged own-load price $p$,
\[
 r_Z(x;p)\ge r_X(x;p)>0.
\]
Erasure of global failure therefore requires a new supported plan and,
necessarily, $D_Y\le\CH_X$. A sufficient condition for survival is
$\min\{D_X,D_Y\}>\CH_X$.

Fix an attained old hull-dual optimizer $p^*$. If every added plan satisfies
\begin{equation}\label{eq:dualsafe}
 c(y)+(p^*)^\top e(y)\ge V_X(p^*),
\end{equation}
then $\CH_Z=\CH_X$. Under this condition, failure is erased exactly when
some added plan attains equality in \eqref{eq:dualsafe} and also satisfies
$p^*\in\partial F(e(y))$. If \eqref{eq:dualsafe} holds with a uniform
margin $\varepsilon>0$, then
\[
 \Delta_Z\ge\min\{\Delta_X,\varepsilon\}>0.
\]
\end{proposition}

These are elementary consequences of the classical dual certificate
(Appendix~\ref{app:expansion}). They provide a useful way to ask whether
a restricted physical mechanism survives genuinely richer alternatives.
With strictly convex supply, supply equality at $p^*$ fixes the old
hull-optimal load; fleet equality then fixes its operating cost. Under
\eqref{eq:dualsafe}, erasure means physically implementing that projected
hull optimum. Without \eqref{eq:dualsafe}, the hull can move, and comparison
with its old value is not enough.

\subsection{Movement additions under unchanged resources}
In the nominal model of Section~\ref{sec:construction}, retain all original
plans and allow additional movements consuming nonnegative energy and
incurring nonnegative extra operating cost. Preserve the two services,
at most two buses with 20 kWh batteries, fee 7, full initial and terminal
inventories, the early 10 and terminal 30 kW power limits, charging windows and access.
There is no energy recovery, inventory transfer or new charging opportunity.

\begin{theorem}[Movement-only stability]\label{thm:movement}
Every feasible added plan is componentwise dominated in purchased-energy
and operating-cost space by an old physical plan. Consequently
\[
 D_Z=97,\qquad\CH_Z=\frac{7591}{80},\qquad
 \Delta_Z=\frac{169}{80}.
\]
\end{theorem}

To see the physical reason, any added plan buys $30+w$ total energy,
$w\ge0$, and at most 10 early kWh. A two-bus plan with early purchase $u$
is dominated by the old two-bus plan $(u,30-u)$ of cost 14. A one-bus
plan must still obtain at least 10 early kWh to complete both services;
the early limit forces $u=10$, and the old one-bus plan $(10,20)$ of
cost 7 dominates it. The same replacement works inside every hull mixture
because $F$ is coordinatewise nondecreasing. The theorem therefore preserves
both values, not merely positivity of their difference.

This is an unchanged-resource result. The original no-movement construction
already includes both possible service partitions and every continuous
charging choice. A new projected point cannot be called an overlooked
feasible duty without explaining which physical assumption admits it.

\subsection{An explicitly enlarged battery choice}
Now change the equipment. Keep buses with 20 kWh batteries and fee 7, and add buses
of capacity $B\in[20,30]$ and fee $q\ge7$. Two of each type are available,
but at most two buses total are used. Every used bus starts and ends at
its own full capacity. All service and charging constraints are unchanged.

Put $a=30-B$. The new one-bus branch buys $x\in[a,10]$ early and
$30-x$ terminally, with operating cost $q$. Every two-bus type assignment
permits the old interval $[0,10]$ and costs at least 14; the old--old
assignment therefore dominates mixed and new--new assignments in projected
economic comparisons. Their physical records and equal-fee ties remain
part of the complete alternative set.

\begin{theorem}[Battery and fee conditions for erasure]\label{thm:battery}
For $20\le B\le30$ and $q\ge7$, let $a=30-B$ and
\[
 A(a)=14+2a-\frac25a^2,\qquad
 C(a)=7+(10-a)\left(\frac25a-2\right).
\]
The enlarged model has a supported physical plan exactly when
\begin{equation}\label{eq:batteryregion}
 \bigl(a\le5\ \hbox{and}\ q=7\bigr)
 \quad\hbox{or}\quad
 \bigl(a>5\ \hbox{and}\ 7\le q\le\min\{A(a),C(a)\}\bigr).
\end{equation}
The inequalities include weak support at ties. At fee $q=7$, this becomes
\begin{equation}\label{eq:batterythreshold}
 B\ge\frac{55-\sqrt{95}}2\simeq22.6266.
\end{equation}
Outside \eqref{eq:batteryregion}, the physical--hull gap is positive.
\end{theorem}

The new branch's physical minimizer is $t=\max\{5,a\}$. If $a\le5$,
its marginal prices at $t=5$ are equal; the new plan minimizes its bill
against the cheaper old one-bus plan only at fee 7. If $a>5$, the two
inequalities in \eqref{eq:batteryregion} compare it respectively with
the old zero-early two-bus plan and old one-bus plan. The proof in
Appendix~\ref{app:battery} checks the entire continuous opportunity set.

The case $B=25$ shows why a new plan's physical cost is not the whole
story. For $q\ge7$,
\begin{align}
 D_Z&=\min\{97,85+q\},\label{eq:B25D}\\
 \CH_Z&=
 \begin{cases}
 85+q-(q-7)^2/20,&7\le q\le21/2,\\
 7591/80,&q\ge21/2.
 \end{cases}\label{eq:B25CH}
\end{align}
At fee 7, the new one-bus load $(5,25)$ has total cost 92 and is supported
at $(5,5)$. At fee 8, its physical cost 93 is below the \emph{old} hull
value 94.8875, but the \emph{new} hull value is 92.95. Its regret is 1 and
the gap remains $1/20$. Thus the necessary condition $D_Y\le\CH_X$
in Proposition~\ref{prop:expansion} is not sufficient.

At fee 12, the old one-bus load $(10,20)$ and new one-bus load $(5,25)$
both have physical cost 97. Their own-price regrets are 13 and 5,
respectively; $\CH_Z=7591/80$. Reporting a single selected planner
optimum would hide this incentive difference. The new equipment can erase
failure, but its presence alone does not do so.
